\documentclass[10pt,a4paper]{amsart}
\usepackage{amsmath,amssymb,amsthm,graphicx}
% Retained minimal preamble for the Stacks Project chapter sites-cohomology.tex (Cohomology on Sites).
% Verbatim excerpt of the chapter preamble: Xy-pic setup (native preamble lines 14-18) and the native
% theorem declarations with the mathematical macros (native preamble lines 154-217). Copied mechanically;
% nothing added, removed or redefined. The wrapper class and the amsmath/amssymb/amsthm/graphicx base
% packages are supplied by the assembler, so the class/externaldocument part of the native preamble is omitted.
% For commutative diagrams we use Xy-pic
\usepackage[all]{xy}

% We use 2cell for 2-commutative diagrams.
\xyoption{2cell}
% Theorem environments.
%
\theoremstyle{plain}
\newtheorem{theorem}[subsection]{Theorem}
\newtheorem{proposition}[subsection]{Proposition}
\newtheorem{lemma}[subsection]{Lemma}

\theoremstyle{definition}
\newtheorem{definition}[subsection]{Definition}
\newtheorem{example}[subsection]{Example}
\newtheorem{exercise}[subsection]{Exercise}
\newtheorem{situation}[subsection]{Situation}

\theoremstyle{remark}
\newtheorem{remark}[subsection]{Remark}
\newtheorem{remarks}[subsection]{Remarks}

\numberwithin{equation}{subsection}

% Macros
%
\def\lim{\mathop{\mathrm{lim}}\nolimits}
\def\colim{\mathop{\mathrm{colim}}\nolimits}
\def\Spec{\mathop{\mathrm{Spec}}}
\def\Hom{\mathop{\mathrm{Hom}}\nolimits}
\def\Ext{\mathop{\mathrm{Ext}}\nolimits}
\def\SheafHom{\mathop{\mathcal{H}\!\mathit{om}}\nolimits}
\def\SheafExt{\mathop{\mathcal{E}\!\mathit{xt}}\nolimits}
\def\Sch{\mathit{Sch}}
\def\Mor{\mathop{\mathrm{Mor}}\nolimits}
\def\Ob{\mathop{\mathrm{Ob}}\nolimits}
\def\Sh{\mathop{\mathit{Sh}}\nolimits}
\def\NL{\mathop{N\!L}\nolimits}
\def\CH{\mathop{\mathrm{CH}}\nolimits}
\def\proetale{{pro\text{-}\acute{e}tale}}
\def\etale{{\acute{e}tale}}
\def\QCoh{\mathit{QCoh}}
\def\Ker{\mathop{\mathrm{Ker}}}
\def\Im{\mathop{\mathrm{Im}}}
\def\Coker{\mathop{\mathrm{Coker}}}
\def\Coim{\mathop{\mathrm{Coim}}}

% Boxtimes
%
\DeclareMathSymbol{\boxtimes}{\mathbin}{AMSa}{"02}

%
% Macros for moduli stacks/spaces
%
\def\QCohstack{\mathcal{QC}\!\mathit{oh}}
\def\Cohstack{\mathcal{C}\!\mathit{oh}}
\def\Spacesstack{\mathcal{S}\!\mathit{paces}}
\def\Quotfunctor{\mathrm{Quot}}
\def\Hilbfunctor{\mathrm{Hilb}}
\def\Curvesstack{\mathcal{C}\!\mathit{urves}}
\def\Polarizedstack{\mathcal{P}\!\mathit{olarized}}
\def\Complexesstack{\mathcal{C}\!\mathit{omplexes}}
% \Pic is the operator that assigns to X its picard group, usage \Pic(X)
% \Picardstack_{X/B} denotes the Picard stack of X over B
% \Picardfunctor_{X/B} denotes the Picard functor of X over B
\def\Pic{\mathop{\mathrm{Pic}}\nolimits}
\def\Picardstack{\mathcal{P}\!\mathit{ic}}
\def\Picardfunctor{\mathrm{Pic}}
\def\Deformationcategory{\mathcal{D}\!\mathit{ef}}
\usepackage{cleveref}
\crefname{originallemma}{original lemma}{original lemmas}
\Crefname{originallemma}{Original lemma}{Original lemmas}
\makeatletter
\newlabel{derived-section-derived-functors}{{05S7}{source}{}{sourcepointer-1}{}}
\newlabel{derived-section-derived-functors@cref}{{[originallemma][0][]05S7}{[1][0][]source}{}{}{}}
\makeatother
\crefname{originallemma}{original lemma}{original lemmas}
\Crefname{originallemma}{Original lemma}{Original lemmas}
\makeatletter
\newlabel{derived-situation-derived-functor}{{05S8}{source}{}{sourcepointer-2}{}}
\newlabel{derived-situation-derived-functor@cref}{{[originallemma][0][]05S8}{[1][0][]source}{}{}{}}
\makeatother
\crefname{originallemma}{original lemma}{original lemmas}
\Crefname{originallemma}{Original lemma}{Original lemmas}
\makeatletter
\newlabel{derived-definition-right-derived-functor-defined}{{05S9}{source}{}{sourcepointer-3}{}}
\newlabel{derived-definition-right-derived-functor-defined@cref}{{[originallemma][0][]05S9}{[1][0][]source}{}{}{}}
\makeatother
\crefname{originallemma}{original lemma}{original lemmas}
\Crefname{originallemma}{Original lemma}{Original lemmas}
\makeatletter
\newlabel{derived-lemma-find-existence-computes}{{06XN}{source}{}{sourcepointer-4}{}}
\newlabel{derived-lemma-find-existence-computes@cref}{{[originallemma][0][]06XN}{[1][0][]source}{}{}{}}
\makeatother
\crefname{originallemma}{original lemma}{original lemmas}
\Crefname{originallemma}{Original lemma}{Original lemmas}
\makeatletter
\newlabel{lemma-K-flat-resolution}{{06YS}{source}{}{sourcepointer-5}{}}
\newlabel{lemma-K-flat-resolution@cref}{{[originallemma][0][]06YS}{[1][0][]source}{}{}{}}
\makeatother
\crefname{originallemma}{original lemma}{original lemmas}
\Crefname{originallemma}{Original lemma}{Original lemmas}
\makeatletter
\newlabel{lemma-pullback-K-flat}{{0G7E}{source}{}{sourcepointer-6}{}}
\newlabel{lemma-pullback-K-flat@cref}{{[originallemma][0][]0G7E}{[1][0][]source}{}{}{}}
\makeatother
\crefname{originallemma}{original lemma}{original lemmas}
\Crefname{originallemma}{Original lemma}{Original lemmas}
\makeatletter
\newlabel{lemma-derived-tor-quasi-isomorphism-other-side}{{06YT}{source}{}{sourcepointer-7}{}}
\newlabel{lemma-derived-tor-quasi-isomorphism-other-side@cref}{{[originallemma][0][]06YT}{[1][0][]source}{}{}{}}
\makeatother
\crefname{originallemma}{original lemma}{original lemmas}
\Crefname{originallemma}{Original lemma}{Original lemmas}
\makeatletter
\newlabel{derived-equation-everywhere}{{05SW}{source}{}{sourcepointer-8}{}}
\newlabel{derived-equation-everywhere@cref}{{[originallemma][0][]05SW}{[1][0][]source}{}{}{}}
\makeatother
\crefname{originallemma}{original lemma}{original lemmas}
\Crefname{originallemma}{Original lemma}{Original lemmas}
\makeatletter
\newlabel{lemma-bounded-flat-K-flat}{{06YQ}{source}{}{sourcepointer-9}{}}
\newlabel{lemma-bounded-flat-K-flat@cref}{{[originallemma][0][]06YQ}{[1][0][]source}{}{}{}}
\makeatother
\title{Existence of the derived pullback agreeing with the naive pullback on K-flat complexes with flat terms}
\author{The Stacks Project Authors}
\date{stacks-project:sites-cohomology.tex:local-snapshot-20261011; GNU FDL 1.2 or later}
\begin{document}
\maketitle
\noindent\textit{Editorial note.} All delivered mathematical text is the author's own, copied from the pinned Stacks Project chapter file sites-cohomology.tex; no statement or proof step was rewritten, reordered or removed.\par\medskip
\noindent\textit{Editorial note.} The proof is the author's complete verification: the hypotheses of the abstract derived-functor result are checked in this situation and the identification of the values is attributed to the corresponding result of this chapter. The whole verification lies inside the counted span.\par\medskip
\noindent\textit{Editorial note.} Every labelled reference occurring inside the counted spans is carried in assembly.external\_references as a non-load-bearing navigation pointer (same-chapter ones at their definition site, cross-chapter ones at the invocation inside the counted span), so no load-bearing mathematics is replaced by a stub. No bibliography citation occurs inside the counted spans.\par\medskip
\noindent\textit{Editorial note.} The checked chapter list for same-name collisions is the set of stacks chapters staged in this repository: chow, cohomology, divisors, etale-cohomology, flat, more-algebra, more-morphisms, morphisms, spaces-more-morphisms, varieties together with this card's chapter sites-cohomology.\par\medskip
\noindent\textit{Editorial note.} The preamble is the chapter's own retained minimal preamble (sources/stacks-sites-cohomology/retained-preamble.tex), a verbatim excerpt of the native preamble: the Xy-pic setup and the native theorem declarations with the mathematical macros.\par\medskip
\noindent\textit{Allowed standard prerequisites (editorial).}\begin{itemize}
\item Derived tensor product and K-flat complexes on a ringed topos. Same chapter, the derived tensor product and the role of K-flat representatives in computing it. Supplies the machinery through which the functor Lf\textasciicircum{}* is constructed.
\item General theory of the derived functor of an additive functor between triangulated categories. Derived Categories, the construction of a left derived functor through a suitable class of complexes and the verification that it is exact and triangulated. Supplies the abstract result which the author applies to the pullback functor f\textasciicircum{}*.
\item Pullback of K-flat complexes with flat terms. Same chapter, the result that the pullback of a K-flat complex with flat terms is again K-flat with flat terms. Supplies the identification Lf\textasciicircum{}*K\textasciicircum{}bullet = f\textasciicircum{}*K\textasciicircum{}bullet for the complexes mentioned in the statement, and the bounded above case through the flatness of the terms.
\end{itemize}\medskip
\ifdefined\hypertarget\hypertarget{sourcepointer-1}{}\fi
\noindent\textit{Source reference.} Original lemma 05S7: sources/stacks-sites-cohomology/pinned-original-source.tex, source lines 3599--3599. Navigation pointer invoked at proof line 3599; the counted span contains the pointer only and the cited result is not reproduced inside it, so the pointer is registered as a non-load-bearing navigation stub in the established form for Stacks material. No bibliography citation is involved.\par\medskip
\ifdefined\hypertarget\hypertarget{sourcepointer-2}{}\fi
\noindent\textit{Source reference.} Original lemma 05S8: sources/stacks-sites-cohomology/pinned-original-source.tex, source lines 3607--3607. Navigation pointer invoked at proof line 3607; the counted span contains the pointer only and the cited result is not reproduced inside it, so the pointer is registered as a non-load-bearing navigation stub in the established form for Stacks material. No bibliography citation is involved.\par\medskip
\ifdefined\hypertarget\hypertarget{sourcepointer-3}{}\fi
\noindent\textit{Source reference.} Original lemma 05S9: sources/stacks-sites-cohomology/pinned-original-source.tex, source lines 3610--3610. Navigation pointer invoked at proof line 3610; the counted span contains the pointer only and the cited result is not reproduced inside it, so the pointer is registered as a non-load-bearing navigation stub in the established form for Stacks material. No bibliography citation is involved.\par\medskip
\ifdefined\hypertarget\hypertarget{sourcepointer-4}{}\fi
\noindent\textit{Source reference.} Original lemma 06XN: sources/stacks-sites-cohomology/pinned-original-source.tex, source lines 3613--3613. Navigation pointer invoked at proof line 3613; the counted span contains the pointer only and the cited result is not reproduced inside it, so the pointer is registered as a non-load-bearing navigation stub in the established form for Stacks material. No bibliography citation is involved.\par\medskip
\ifdefined\hypertarget\hypertarget{sourcepointer-5}{}\fi
\noindent\textit{Source reference.} Original lemma 06YS: sources/stacks-sites-cohomology/pinned-original-source.tex, source lines 3207--3207. Navigation pointer invoked at proof line 3616; the counted span contains the pointer only and the cited result is not reproduced inside it, so the pointer is registered as a non-load-bearing navigation stub in the established form for Stacks material. No bibliography citation is involved.\par\medskip
\ifdefined\hypertarget\hypertarget{sourcepointer-6}{}\fi
\noindent\textit{Source reference.} Original lemma 0G7E: sources/stacks-sites-cohomology/pinned-original-source.tex, source lines 3508--3508. Navigation pointer invoked at proof line 3632; the counted span contains the pointer only and the cited result is not reproduced inside it, so the pointer is registered as a non-load-bearing navigation stub in the established form for Stacks material. No bibliography citation is involved.\par\medskip
\ifdefined\hypertarget\hypertarget{sourcepointer-7}{}\fi
\noindent\textit{Source reference.} Original lemma 06YT: sources/stacks-sites-cohomology/pinned-original-source.tex, source lines 3233--3233. Navigation pointer invoked at proof line 3636; the counted span contains the pointer only and the cited result is not reproduced inside it, so the pointer is registered as a non-load-bearing navigation stub in the established form for Stacks material. No bibliography citation is involved.\par\medskip
\ifdefined\hypertarget\hypertarget{sourcepointer-8}{}\fi
\noindent\textit{Source reference.} Original lemma 05SW: sources/stacks-sites-cohomology/pinned-original-source.tex, source lines 3646--3646. Navigation pointer invoked at proof line 3646; the counted span contains the pointer only and the cited result is not reproduced inside it, so the pointer is registered as a non-load-bearing navigation stub in the established form for Stacks material. No bibliography citation is involved.\par\medskip
\ifdefined\hypertarget\hypertarget{sourcepointer-9}{}\fi
\noindent\textit{Source reference.} Original lemma 06YQ: sources/stacks-sites-cohomology/pinned-original-source.tex, source lines 3120--3120. Navigation pointer invoked at proof line 3539; the counted span contains the pointer only and the cited result is not reproduced inside it, so the pointer is registered as a non-load-bearing navigation stub in the established form for Stacks material. No bibliography citation is involved.\par\medskip
\begin{lemma}
\label{lemma-derived-base-change}
Let $f : (\Sh(\mathcal{C}), \mathcal{O}) \to (\Sh(\mathcal{C}'), \mathcal{O}')$
be a morphism of ringed topoi. There exists an exact functor
$$
Lf^* : D(\mathcal{O}') \longrightarrow D(\mathcal{O})
$$
of triangulated categories so that
$Lf^*\mathcal{K}^\bullet = f^*\mathcal{K}^\bullet$ for any
K-flat complex $\mathcal{K}^\bullet$ with flat terms and
in particular for any bounded above complex of flat $\mathcal{O}'$-modules.
\end{lemma}

\begin{proof}
To see this we use the general theory developed in
Derived Categories, Section \ref{derived-section-derived-functors}.
Set $\mathcal{D} = K(\mathcal{O}')$ and $\mathcal{D}' = D(\mathcal{O})$.
Let us write $F : \mathcal{D} \to \mathcal{D}'$ the exact functor
of triangulated categories defined by the rule
$F(\mathcal{G}^\bullet) = f^*\mathcal{G}^\bullet$.
We let $S$ be the set of quasi-isomorphisms in
$\mathcal{D} = K(\mathcal{O}')$.
This gives a situation as in
Derived Categories, Situation \ref{derived-situation-derived-functor}
so that
Derived Categories, Definition
\ref{derived-definition-right-derived-functor-defined}
applies. We claim that $LF$ is everywhere defined.
This follows from
Derived Categories, Lemma \ref{derived-lemma-find-existence-computes}
with $\mathcal{P} \subset \Ob(\mathcal{D})$ the collection
of K-flat complexes $\mathcal{K}^\bullet$ with flat terms.
Namely, (1) follows from Lemma \ref{lemma-K-flat-resolution}
and to see (2) we have to show that for a quasi-isomorphism
$\mathcal{K}_1^\bullet  \to \mathcal{K}_2^\bullet$ between
elements of $\mathcal{P}$ the map
$f^*\mathcal{K}_1^\bullet  \to f^*\mathcal{K}_2^\bullet$ is a
quasi-isomorphism. To see this write this as
$$
f^{-1}\mathcal{K}_1^\bullet \otimes_{f^{-1}\mathcal{O}'} \mathcal{O}
\longrightarrow
f^{-1}\mathcal{K}_2^\bullet \otimes_{f^{-1}\mathcal{O}'} \mathcal{O}
$$
The functor $f^{-1}$ is exact, hence the map
$f^{-1}\mathcal{K}_1^\bullet  \to f^{-1}\mathcal{K}_2^\bullet$ is a
quasi-isomorphism. The complexes
$f^{-1}\mathcal{K}_1^\bullet$ and $f^{-1}\mathcal{K}_2^\bullet$
are K-flat complexes of $f^{-1}\mathcal{O}'$-modules by
Lemma \ref{lemma-pullback-K-flat}
because we can consider the morphism of ringed topoi
$(\Sh(\mathcal{C}), f^{-1}\mathcal{O}') \to
(\Sh(\mathcal{C}'), \mathcal{O}')$. Hence
Lemma \ref{lemma-derived-tor-quasi-isomorphism-other-side}
guarantees that the displayed map is a quasi-isomorphism.
Thus we obtain a derived functor
$$
LF :
D(\mathcal{O}') = S^{-1}\mathcal{D}
\longrightarrow
\mathcal{D}' = D(\mathcal{O})
$$
see
Derived Categories, Equation (\ref{derived-equation-everywhere}).
Finally,
Derived Categories, Lemma \ref{derived-lemma-find-existence-computes}
also guarantees that
$LF(\mathcal{K}^\bullet) = F(\mathcal{K}^\bullet) = f^*\mathcal{K}^\bullet$
when $\mathcal{K}^\bullet$ is in $\mathcal{P}$.
The proof is finished by observing that
bounded above complexes of flat modules are in $\mathcal{P}$
by Lemma \ref{lemma-bounded-flat-K-flat}.
\end{proof}

\end{document}
